\chapter{Калибровочное сокращение неподвижной алгебры и двухсекторный остаток}
\label{ch:v8-markov-fixed-algebra-selector-gate}

\section{Проверяемая щель}

Связывающая квантово-марковская полугруппа имеет неподвижную алгебру
размерности $41$. Выбирать внутри неё максимальную коммутативную
подалгебру вручную запрещено. Поэтому проверим только операторы, уже
присутствующие на физическом концевом носителе: восемь генераторов
$SU(3)$, три генератора $SU(2)$, гиперзаряд и связывающую инцидентность $A_0$.

Для каждого самосопряжённого генератора $G$ добавляется форма Дирихле
\begin{equation}
 \mathcal L_G=-\frac12\operatorname{ad}_G^2.
\end{equation}
Нас интересует не относительная скорость этих процессов, а общее ядро.
Для любых строго положительных весов оно равно пересечению коммутантов и
потому не зависит от нормировок калибровочных множителей.

\section{Последовательное снятие кратностей}

На концевой алгебре $M_{11}\oplus M_{10}$ аудит даёт
\begin{center}
\begin{tabular}{lc}
\toprule
Набор генераторов & Размерность неподвижной алгебры\\
\midrule
только связывающую инцидентность & $41$\\
$A_0+U(1)_Y$ & $21$\\
$A_0+SU(2)$ & $10$\\
$A_0+SU(3)$ & $9$\\
$A_0+SU(3)+U(1)_Y$ & $5$\\
полный $A_0+SU(3)+SU(2)+U(1)_Y$ & $2$\\
\bottomrule
\end{tabular}
\end{center}

Два оставшихся направления порождены ортогональными центральными
проекторами
\begin{equation}
 P_q=P_{Q_L}\oplus P_{u_R}\oplus P_{d_R},
 \qquad \rank P_q=6+3+3=12,
 \label{eq:v8-markov-quark-projector}
\end{equation}
и
\begin{equation}
 P_\ell=I-P_q,
 \qquad \rank P_\ell=9.
 \label{eq:v8-markov-lepton-projector}
\end{equation}
Второй сектор содержит $L_L,X_L,Y_L,e_R,X_R,Y_R$. Оба проектора точно
аннулируются полным генератором в машинной точности.

\section{Устойчивость по метрике Дирихле}

Независимые веса связывающей, $SU(3)$-, $SU(2)$- и $U(1)$-форм сканировались
в диапазоне от $10^{-4}$ до $10^4$. Во всех шестидесяти четырёх тестах
размерность общего ядра оставалась равной двум. Следовательно, результат
не использует невыведенное отношение калибровочных констант.

Градуировка концевых секторов скалярна на каждом углу, поэтому её двойной коммутатор
на $M_{11}\oplus M_{10}$ тождественно равен нулю. Аффинный семейный фактор
действует здесь как единица и также не снимает двухсекторный остаток.

\begin{theorem}[Канонический двухсекторный остаток]
Совместная форма Дирихле физической связывающей инцидентности и полного
представления калибровочной алгебры сокращает 41-мерную неподвижную алгебру
ровно до
\begin{equation}
 \mathcal Z_{\rm fix}=\mathbb CP_q\oplus\mathbb CP_\ell\simeq\mathbb C^2.
\end{equation}
Это каноническая коммутативная секторная алгебра, независимая от всех
строго положительных относительных весов. Одновременно она доказывает, что
текущая инцидентность не содержит перехода между цветным кварковым сектором и
бесцветным лептонный/вектороподобный-сектором.
\end{theorem}

\begin{remark}
Полученная $\mathbb C^2$ не является пространством отдельных частиц или
точек пространства-времени. Она различает только две связные компоненты
внутреннего процесса. Положительная часть результата --- вывод первого
классического секторного наблюдаемого. Отрицательная --- точная фиксация
недостающего межкомпонентного перехода.
\end{remark}

\begin{nogocriterion}[Запрет смешивания центральных секторов вручную]
Нельзя добавлять произвольный блок между $P_q$ и $P_\ell$ лишь для получения
примитивной полугруппы. Новый переход обязан быть совместимым с калибровочной симметрией,
возникать из полного корреляционного ядра или составной кривизны и проходить
старые запреты на цветной вакуум и невыведенный соединитель Мориты.
\end{nogocriterion}

\begin{protocol}
\begin{enumerate}
 \item исходная неподвижная размерность: \textbf{$41$};
 \item после полного калибровочный набора: \textbf{$2$};
 \item каноническая коммутативная алгебра: \textbf{$\mathbb C^2$};
 \item ранги центральных проекторов: \textbf{$12$ и $9$};
 \item устойчивость при положительных весах: \textbf{пройдена};
 \item градуировка концевых секторов снимает остаток: \textbf{нет};
 \item аффинный фактор снимает остаток: \textbf{нет};
 \item переход кварк--лептон: \textbf{отсутствует};
 \item единая примитивная полугруппа: \textbf{не получена};
 \item двухсекторная граница: \textbf{строго зафиксирована}.
\end{enumerate}
\end{protocol}

Машинный сертификат сохранён в
\path{s2t_v8_markov_fixed_algebra_selector_gate_results.json}.